\Qhead{4.4}{Оптические компрессоры и получение фемтосекундных импульсов}{Оптические компрессоры и получение фемтосекундных импульсов.}
\Subfile{4.4.1}{sections/04_generation/q04/s01}

\recap{Короткий импульс требует широкой полосы и согласованной спектральной фазы. Компрессор компенсирует частотную задержку; CPA разделяет увеличение энергии и пиковой мощности.}
\board{Спектр с квадратичной фазой; растяжитель–усилитель–компрессор; произведение ширин гауссова импульса.}
\oral{Можно ли определить длительность чирпованного импульса по одному измерению его спектральной интенсивности?}
\bait{Можно оставить высшие порядки фазы: компенсация только квадратичного чирпа не обеспечивает предельного сжатия широкой полосы.}
\sources{Kärtner~\cite{kartnerfull}, главы о дисперсии, коротких импульсах и их измерении; связь со вторым приближением дисперсии — вопрос~\Qref{2.2}.}
